%% kapa: exam=2 round=31 year=2020 period=1 subject=appraisal_practice subject_ko=감정평가실무 seq=1 number=1 points=40
%% source: Q-Net 감정평가사 기출문제 (cst00309) artlSeq=5207673 file=raw/1008/5207673/1교시_감정평가실무.pdf
%% conversion: kordoc 4.24.51 -> tools/gen2.py (auto); review=proofread-gpt-astra; flags=none
\begin{longquestion}{【문제 1】}
주식회사A는 주식회사B를 인수합병하는 프로젝트에서 주식회사B의 영업권 가치를 파악하기 위해 감정평가사甲에게 감정평가를 의뢰하였다. 관련 법규 및 이론을 참작하고 제시된 자료를 활용하여 다음의 물음에 답하시오. (단, 기준시점은 2020. 09. 19임) (40점) \par
\medskip\noindent 물음 1) 주식회사B 소유 부동산의 공정가치를 3방식을 적용하여 감정평가하고, 시산가액 조정을 통해 결정하시오. (30점) \par
\medskip\noindent 물음 2) 「감정평가에 관한 규칙」, 「감정평가 실무기준」에 의거하여 영업권 가치를 감정평가 하시오. (10점) \par
< 자료 1 > 공통사항 \par
1. 단가는 유효숫자 셋째자리까지 표시하되, 넷째자리 이하는 절사함 \par
2. 시산가액과 총액은 백만원 단위까지 표시하되, 십만원 단위이하는 절사함 \par
< 자료 2 > 주식회사B 소유 부동산 현황 \par
1. 평가대상(집합건물) 물건 \par
{\small\centering\renewcommand{\arraystretch}{1.25}
\resizebox{\ifdim\width>\linewidth\linewidth\else\width\fi}{!}{
\begin{tabular}{|c|c|c|c|c|c|}
\hline
\multicolumn{2}{|c|}{\makecell{소재지}} & \multicolumn{4}{|c|}{\makecell{S시 K구 J동 100-1번지 (S시 K구 OO로 5)}} \\
\hline
\makecell{층} & \makecell{호수} & \makecell{용도} & \makecell{전유면적(m\textsuperscript{2})} & \makecell{공용면적(m\textsuperscript{2})} & \makecell{계약면적(m\textsuperscript{2})} \\
\hline
\makecell{지하1층} & \makecell{B101} & \makecell{근린생활시설} & \makecell{1,200} & \makecell{630} & \makecell{1,830} \\
\hline
\makecell{지상1층} & \makecell{101} & \makecell{근린생활시설} & \makecell{950} & \makecell{500} & \makecell{1,450} \\
\hline
\makecell{지상2층} & \makecell{201} & \makecell{업무시설} & \makecell{1,200} & \makecell{630} & \makecell{1,830} \\
\hline
\makecell{지상3층} & \makecell{301} & \makecell{업무시설} & \makecell{1,200} & \makecell{630} & \makecell{1,830} \\
\hline
\makecell{지상4층} & \makecell{401} & \makecell{업무시설} & \makecell{1,200} & \makecell{630} & \makecell{1,830} \\
\hline
\makecell{지상5층} & \makecell{501} & \makecell{업무시설} & \makecell{1,000} & \makecell{520} & \makecell{1,520} \\
\hline
\end{tabular}}\par}

2. 대상 토지 현황 \par
{\small\centering\renewcommand{\arraystretch}{1.25}
\resizebox{\ifdim\width>\linewidth\linewidth\else\width\fi}{!}{
\begin{tabular}{|p{0.23\linewidth}|p{0.70\linewidth}|}
\hline
\makecell{위치 및 주위환경} & 본 건은 S시 K구 J동 소재 S시청 동측 인근에 위치하고 S시 도심지역에 속하며, 본 건 주위는 각종 업무용 빌딩, 근린생활시설, 공공청사 등이 소재하는 업무지대임 \\
\hline
\makecell{교통상황} & 본 건까지 차량 진출입 가능하고, 인근에 시내버스 정류장이 소재하며, 지하철 1호선 “○○역”, 2호선 “○○역”이 소재하여 제반 교통상황은 양호함 \\
\hline
\makecell{지목 / 면적} & 대 / 1,800m\textsuperscript{2} \\
\hline
\makecell{형상, 지세 및\\이용상황} & 인접지와 등고 평탄한 가장형 토지로서, 업무용 건부지로 이용중임 \\
\hline
\makecell{접면도로} & 서측으로 폭 약 40m, 남측으로 폭 약 10m 포장도로에 각각 접함 \\
\hline
\makecell{토지이용계획 및\\공법상 제한사항} & 일반상업지역 \\
\hline
\end{tabular}}\par}

3. 대상 건물 현황 \par
{\small\centering\renewcommand{\arraystretch}{1.25}
\resizebox{\ifdim\width>\linewidth\linewidth\else\width\fi}{!}{
\begin{tabular}{|p{0.23\linewidth}|p{0.70\linewidth}|}
\hline
\makecell{건물명} & B빌딩 \\
\hline
\makecell{주용도} & 업무시설, 근린생활시설 \\
\hline
\makecell{건축규모} & 지하2층 / 지상5층 \\
\hline
\makecell{연면적} & 10,290m\textsuperscript{2} \\
\hline
\makecell{호수 / 사용승인일} & 6개호 / 2000.12.31 \\
\hline
\makecell{구조 및 지붕} & 철골철근콘크리트조 슬래브지붕 \\
\hline
\makecell{마감재} & 외벽: 화강석 및 복합판넬 마감 등 내벽: 몰탈 위 페인팅, 타일붙임 및 내부 인테리어 마감 등 창호: 강화 유리창 마감 등 \\
\hline
\makecell{층별 용도\\(임대현황 등)} & 지하2층(1,695m\textsuperscript{2}): 주차장, 기계실 지하1층(1,695m\textsuperscript{2}): 근린생활시설(임대: 음식점), 주차장 지상1층(1,150m\textsuperscript{2}): 근린생활시설(임대: 카페) 지상2층(1,500m\textsuperscript{2}): 업무시설(임대: W법무법인) 지상3층(1,500m\textsuperscript{2}): 업무시설(자가사용: 주식회사B) 지상4층(1,500m\textsuperscript{2}): 업무시설(자가사용: 주식회사B) 지상5층(1,250m\textsuperscript{2}): 업무시설(자가사용: 주식회사B) \\
\hline
\end{tabular}}\par}

< 자료 3 > 비교방식 참고자료 \par
1. 인근지역 집합건물 거래사례(소재지: S시 K구 J동) \par
{\small\centering\renewcommand{\arraystretch}{1.25}
\resizebox{\ifdim\width>\linewidth\linewidth\else\width\fi}{!}{
\begin{tabular}{|c|c|c|c|c|c|c|c|c|c|}
\hline
\makecell{기호} & \makecell{지번} & \makecell{층} & \makecell{호수} & \makecell{용도} & \makecell{전유면적\\(m\textsuperscript{2})} & \makecell{거래가격\\(천원)} & \makecell{거래시점} & \makecell{용도지역} & \makecell{사용\\승인일} \\
\hline
\makecell{1} & \makecell{50} & \makecell{3} & \makecell{301} & \makecell{업무시설} & \makecell{1,000} & \makecell{5,000,000} & \makecell{2020.02.15} & \makecell{준주거} & \makecell{2005.10.31} \\
\hline
\makecell{2} & \makecell{70} & \makecell{1} & \makecell{101} & \makecell{근린생활시설} & \makecell{750} & \makecell{9,750,000} & \makecell{2020.03.20} & \makecell{일반상업} & \makecell{2000.09.15} \\
\hline
\makecell{3} & \makecell{90} & \makecell{2} & \makecell{205} & \makecell{근린생활시설} & \makecell{100} & \makecell{900,000} & \makecell{2019.08.20} & \makecell{일반상업} & \makecell{2003.06.10} \\
\hline
\makecell{4} & \makecell{95} & \makecell{6} & \makecell{601} & \makecell{업무시설} & \makecell{1,050} & \makecell{6,825,000} & \makecell{2020.01.20} & \makecell{일반상업} & \makecell{2002.10.20} \\
\hline
\end{tabular}}\par}

※ 상기 거래사례는 정상적인 거래임 \par
2. 지역별 자본수익률 \par
\medskip\noindent (1) 집합상가(S시 도심지역) \par
(단위: \%) \par
{\small\centering\renewcommand{\arraystretch}{1.25}
\resizebox{\ifdim\width>\linewidth\linewidth\else\width\fi}{!}{
\begin{tabular}{|c|c|c|c|c|}
\hline
\makecell{구분} & \makecell{1분기} & \makecell{2분기} & \makecell{3분기} & \makecell{4분기} \\
\hline
\makecell{2019년} & \makecell{0.42} & \makecell{0.75} & \makecell{0.92} & \makecell{0.78} \\
\hline
\makecell{2020년} & \makecell{0.35} & \makecell{0.32} & \makecell{-} & \makecell{-} \\
\hline
\end{tabular}}\par}

\medskip\noindent (2) 오피스(S시 도심지역) \par
(단위: \%) \par
{\small\centering\renewcommand{\arraystretch}{1.25}
\resizebox{\ifdim\width>\linewidth\linewidth\else\width\fi}{!}{
\begin{tabular}{|c|c|c|c|c|}
\hline
\makecell{구분} & \makecell{1분기} & \makecell{2분기} & \makecell{3분기} & \makecell{4분기} \\
\hline
\makecell{2019년} & \makecell{0.57} & \makecell{0.92} & \makecell{1.13} & \makecell{0.93} \\
\hline
\makecell{2020년} & \makecell{0.54} & \makecell{0.48} & \makecell{-} & \makecell{-} \\
\hline
\end{tabular}}\par}

\medskip\noindent (3) 2020년 3분기 이후 자본수익률은 기준시점 현재 발표되지 않아 2020년 2분기 자본수익률을 연장 추정하여 적용하되, 소수점 넷째자리에서 반올림하여 셋째자리까지 표시함 \par
3. 가치형성요인 비교 참고자료 \par
\medskip\noindent (1) 선정된 거래사례와 대상 물건은 층별효용을 제외한 가치형성요인은 동일함 \par
\medskip\noindent (2) 인근지역 내 대상 물건과 이용상황이 유사한 부동산의 임대료 수준, 평가사례 등을 종합적으로 고려한 결과 다음과 같은 층별효용비를 도출하였음(3방식 공통 적용) \par
{\small\centering\renewcommand{\arraystretch}{1.25}
\resizebox{\ifdim\width>\linewidth\linewidth\else\width\fi}{!}{
\begin{tabular}{|c|c|c|}
\hline
\makecell{층} & \makecell{용도} & \makecell{효용비} \\
\hline
\makecell{지하1층} & \makecell{근린생활시설} & \makecell{35} \\
\hline
\makecell{지상1층} & \makecell{근린생활시설} & \makecell{100} \\
\hline
\makecell{지상2층～최상층} & \makecell{업무시설} & \makecell{50} \\
\hline
\end{tabular}}\par}

< 자료 4 > 원가방식 참고자료 \par
1. 표준지 공시지가(공시기준일: 2020. 01. 01) \par
{\small\centering\renewcommand{\arraystretch}{1.25}
\resizebox{\ifdim\width>\linewidth\linewidth\else\width\fi}{!}{
\begin{tabular}{|c|c|c|c|c|c|c|c|c|c|}
\hline
\makecell{기호} & \makecell{소재지} & \makecell{면적\\(m\textsuperscript{2})} & \makecell{지목} & \makecell{이용\\상황} & \makecell{용도지역} & \makecell{도로\\접면} & \makecell{형상\\지세} & \makecell{공시지가\\(원/m\textsuperscript{2})} & \makecell{비고} \\
\hline
\makecell{A} & \makecell{S시 K구\\J동 97-1} & \makecell{1,332} & \makecell{대} & \makecell{상업용} & \makecell{일반상업} & \makecell{광대\\세각} & \makecell{가장형\\평지} & \makecell{14,500,000} & \makecell{도로\\15\%저촉} \\
\hline
\makecell{B} & \makecell{S시 K구\\J동 98-1} & \makecell{1,665} & \makecell{대} & \makecell{업무용} & \makecell{준주거} & \makecell{중로\\한면} & \makecell{가장형\\평지} & \makecell{10,200,000} & \makecell{-} \\
\hline
\end{tabular}}\par}

※ 표준지 기호A의 도시계획시설(도로)은 장기미집행 도시계획시설로서 2020년 7월 1일자로 해제되었음 \par
2. 인근지역 거래사례(소재지: S시 K구 J동) \par
{\small\centering\renewcommand{\arraystretch}{1.25}
\resizebox{\ifdim\width>\linewidth\linewidth\else\width\fi}{!}{
\begin{tabular}{|c|c|c|c|c|c|c|c|c|}
\hline
\multirow{2}{*}{\makecell{기호}} & \multirow{2}{*}{\makecell{지번}} & \multirow{2}{*}{\makecell{용도지역}} & \multirow{2}{*}{\makecell{지목\\이용상황}} & \multirow{2}{*}{\makecell{형상\\도로조건}} & \makecell{토지면적(m\textsuperscript{2})} & \multirow{2}{*}{\makecell{거래가격\\(천원)}} & \multirow{2}{*}{\makecell{거래시점}} & \multirow{2}{*}{\makecell{사용\\승인일}} \\
\cline{6-6}
 &  &  &  &  & \makecell{건물면적(m\textsuperscript{2})} &  &  &  \\
\hline
\multirow{3}{*}{\makecell{a}} & \multirow{2}{*}{\makecell{90-1}} & \multirow{2}{*}{\makecell{일반상업}} & \multirow{2}{*}{\makecell{대\\상업용}} & \multirow{2}{*}{\makecell{세장형\\광대소각}} & \makecell{215} & \multirow{2}{*}{\makecell{4,100,000}} & \multirow{2}{*}{\makecell{2020.03.31}} & \multirow{2}{*}{\makecell{-}} \\
\cline{6-6}
 &  &  &  &  & \makecell{-} &  &  &  \\
\cline{2-9}
 & \multicolumn{8}{|c|}{\makecell{- 매매계약서상 특약사항: 위 매매대금 중 200,000,000원은 지상에 소재하는\\무허가 건축물의 거래대금인 것으로 양자 합의함}} \\
\hline
\multirow{3}{*}{\makecell{b}} & \multirow{2}{*}{\makecell{93-2}} & \multirow{2}{*}{\makecell{일반상업}} & \multirow{2}{*}{\makecell{대\\상업용}} & \multirow{2}{*}{\makecell{가장형\\광대한면}} & \makecell{520} & \multirow{2}{*}{\makecell{12,500,000}} & \multirow{2}{*}{\makecell{2019.03.31}} & \multirow{2}{*}{\makecell{1995.06.20}} \\
\cline{6-6}
 &  &  &  &  & \makecell{3,250} &  &  &  \\
\cline{2-9}
 & \multicolumn{8}{|c|}{\makecell{- 기준시점 건물 재조달원가: 1,350,000원/m\textsuperscript{2}\\- 거래시점 건물 재조달원가: 1,300,000원/m\textsuperscript{2}\\- 구조 및 내용년수: 철골철근콘크리트구조 슬래브지붕, 55년}} \\
\hline
\end{tabular}}\par}

※ 기타 거래조건은 통상적인 것으로 전제함 \par
3. 인근지역 평가사례(소재지: S시 K구 J동) \par
{\small\centering\renewcommand{\arraystretch}{1.25}
\resizebox{\ifdim\width>\linewidth\linewidth\else\width\fi}{!}{
\begin{tabular}{|c|c|c|c|c|c|c|c|c|c|c|}
\hline
\makecell{기호} & \makecell{지번} & \makecell{지목} & \makecell{면적\\(m\textsuperscript{2})} & \makecell{이용상황} & \makecell{용도지역} & \makecell{도로\\교통} & \makecell{형상\\지세} & \makecell{토지단가\\(원/m\textsuperscript{2})} & \makecell{기준시점} & \makecell{평가목적} \\
\hline
\makecell{ㄱ} & \makecell{101} & \makecell{대} & \makecell{320} & \makecell{상업용} & \makecell{일반상업} & \makecell{중로한면} & \makecell{가장형\\평지} & \makecell{15,500,000} & \makecell{2020.08.01} & \makecell{일반거래} \\
\hline
\makecell{ㄴ} & \makecell{102} & \makecell{대} & \makecell{1,500} & \makecell{업무용} & \makecell{일반상업} & \makecell{광대소각} & \makecell{정방형\\평지} & \makecell{19,800,000} & \makecell{2020.01.01} & \makecell{자산\\재평가} \\
\hline
\makecell{ㄷ} & \makecell{103} & \makecell{대} & \makecell{1,750} & \makecell{업무용} & \makecell{일반상업} & \makecell{광대소각} & \makecell{가장형\\평지} & \makecell{21,500,000} & \makecell{2020.05.01} & \makecell{담보} \\
\hline
\end{tabular}}\par}

4. 지가변동률(S시 K구) \par
{\small\centering\renewcommand{\arraystretch}{1.25}
\resizebox{\ifdim\width>\linewidth\linewidth\else\width\fi}{!}{
\begin{tabular}{|c|c|c|c|}
\hline
\multirow{2}{*}{\makecell{기간}} & \multicolumn{2}{|c|}{\makecell{지가변동률(\%)}} & \multirow{2}{*}{\makecell{비고}} \\
\cline{2-3}
 & \makecell{상업지역} & \makecell{주거지역} &  \\
\hline
\makecell{2019.03.31～2019.12.31} & \makecell{1.745} & \makecell{2.345} & \makecell{누계치} \\
\hline
\makecell{2020.01.01～2020.08.31} & \makecell{1.323} & \makecell{1.565} & \makecell{누계치} \\
\hline
\makecell{2020.08.01～2020.08.31} & \makecell{0.254} & \makecell{0.367} & \makecell{2020.08. 당월} \\
\hline
\end{tabular}}\par}

※ 2020.09. 이후 지가변동률은 기준시점 현재 고시되지 않아 2020.08. 지가변동률을 연장 추정하여 적용하되, 소수점 넷째자리에서 반올림하여 셋째자리까지 표시함 \par
5. 지역요인 비교: 대상 부동산, 표준지, 거래사례, 평가전례 모두 인근지역 내에 소재하여 지역요인은 동일함 \par
6. 개별요인 비교 참고자료 \par
\medskip\noindent (1) 가로조건 \par
{\small\centering\renewcommand{\arraystretch}{1.25}
\resizebox{\ifdim\width>\linewidth\linewidth\else\width\fi}{!}{
\begin{tabular}{|c|c|c|c|c|c|}
\hline
\makecell{구분} & \makecell{광대한면} & \makecell{중로한면} & \makecell{소로한면} & \makecell{세로(가)} & \makecell{세로(불)} \\
\hline
\makecell{격차율} & \makecell{1.00} & \makecell{0.85} & \makecell{0.80} & \makecell{0.75} & \makecell{0.65} \\
\hline
\end{tabular}}\par}

\medskip\noindent (2) 획지조건 \par
1) 형상 \par
{\small\centering\renewcommand{\arraystretch}{1.25}
\resizebox{\ifdim\width>\linewidth\linewidth\else\width\fi}{!}{
\begin{tabular}{|c|c|c|c|c|c|}
\hline
\makecell{구분} & \makecell{정방형} & \makecell{가장형} & \makecell{세장형} & \makecell{사다리} & \makecell{부정형} \\
\hline
\makecell{격차율} & \makecell{1.00} & \makecell{1.02} & \makecell{0.97} & \makecell{0.95} & \makecell{0.90} \\
\hline
\end{tabular}}\par}

2) 접면도로상태 \par
{\small\centering\renewcommand{\arraystretch}{1.25}
\resizebox{\ifdim\width>\linewidth\linewidth\else\width\fi}{!}{
\begin{tabular}{|c|c|c|c|}
\hline
\makecell{구분} & \makecell{한면} & \makecell{세각} & \makecell{소각} \\
\hline
\makecell{격차율} & \makecell{1.00} & \makecell{1.03} & \makecell{1.05} \\
\hline
\end{tabular}}\par}

\medskip\noindent (3) 행정적조건 \par
1) 용도지역 \par
{\small\centering\renewcommand{\arraystretch}{1.25}
\resizebox{\ifdim\width>\linewidth\linewidth\else\width\fi}{!}{
\begin{tabular}{|c|c|c|}
\hline
\makecell{구분} & \makecell{일반상업} & \makecell{준주거} \\
\hline
\makecell{격차율} & \makecell{1.00} & \makecell{0.85} \\
\hline
\end{tabular}}\par}

2) 도시계획시설 \par
{\small\centering\renewcommand{\arraystretch}{1.25}
\resizebox{\ifdim\width>\linewidth\linewidth\else\width\fi}{!}{
\begin{tabular}{|c|c|c|}
\hline
\makecell{구분} & \makecell{일반} & \makecell{도로} \\
\hline
\makecell{격차율} & \makecell{1.00} & \makecell{0.85} \\
\hline
\end{tabular}}\par}

\medskip\noindent (4) 상기에서 제시되지 않은 개별요인은 대상 부동산, 표준지, 거래사례, 평가전례 모두 동일함 \par
\medskip\noindent (5) 항목, 세항목간 및 조건단위 격차율은 상승식으로 산정하되, 조건단위 비교치는 소수점 셋째자리에서 반올림하여 둘째자리까지 표시함 \par
7. 그 밖의 요인 보정: 평가사례 중 가장 적정하다고 판단되는 하나를 선택하되, 그 밖의 요인 보정방법은 표준지 기준 산정방식을 적용함 \par
8. 재조달원가(기준시점) \par
{\small\centering\renewcommand{\arraystretch}{1.25}
\resizebox{\ifdim\width>\linewidth\linewidth\else\width\fi}{!}{
\begin{tabular}{|c|c|c|c|c|}
\hline
\makecell{용도} & \makecell{구조} & \makecell{급수} & \makecell{단가(원/m\textsuperscript{2})} & \makecell{내용 년수} \\
\hline
\makecell{근린생활시설} & \makecell{철골철근콘크리트조\\슬래브지붕 (6층이하)} & \makecell{1급} & \makecell{1,300,000} & \makecell{55 (50-60)} \\
\hline
\makecell{사무실} & \makecell{철골철근콘크리트조\\슬래브지붕 (6층이하)} & \makecell{1급} & \makecell{1,500,000} & \makecell{55 (50-60)} \\
\hline
\end{tabular}}\par}

※ 상기 자료는 지상층 기준 단가이고, 지하층은 용도와 상관없이 지상1층 단가의 70\%를 적용함 \par
9. 기타사항 \par
\medskip\noindent (1) 토지는 공시지가기준법과 거래사례비교법에 의한 시산가액을 산정하여 시산가액 조정함 \par
\medskip\noindent (2) 공시지가기준법과 거래사례비교법에서 비교표준지 및 사례선정 시 가장 적정하다고 판단되는 하나를 선정하되 그 사유를 적시할 것 \par
\medskip\noindent (3) 요인비교치는 소수점 넷째자리에서 반올림하여 셋째자리까지 표시함 \par
\medskip\noindent (4) 층별효용비율은 백분율 기준 소수점 둘째자리에서 반올림하여 첫째자리까지 표시함 \par
< 자료 5 > 수익방식 참고자료 \par
1. 대상 부동산의 현황 임대료 \par
{\small\centering\renewcommand{\arraystretch}{1.25}
\resizebox{\ifdim\width>\linewidth\linewidth\else\width\fi}{!}{
\begin{tabular}{|c|c|c|c|c|}
\hline
\makecell{호수} & \makecell{용도} & \makecell{보증금(원/m\textsuperscript{2})} & \makecell{월임대료(원/m\textsuperscript{2})} & \makecell{월관리비(원/m\textsuperscript{2})} \\
\hline
\makecell{B101} & \makecell{근린생활시설} & \makecell{150,000} & \makecell{15,000} & \makecell{3,000} \\
\hline
\makecell{101} & \makecell{근린생활시설} & \makecell{450,000} & \makecell{45,000} & \makecell{5,000} \\
\hline
\makecell{201} & \makecell{업무시설} & \makecell{190,000} & \makecell{19,000} & \makecell{5,000} \\
\hline
\makecell{301～501} & \makecell{업무시설} & \multicolumn{3}{|c|}{\makecell{자가사용}} \\
\hline
\end{tabular}}\par}

※ 상기 임대내역은 계약면적 기준임 \par
※ 현황 임대료는 인근지역 내 평균적인 임대료 수준과 유사한 것으로 조사됨 \par
2. 인근지역 유사 부동산의 표준적 공실 및 대손충당금 비율은 가능총소득(PGI)의 10\%임 \par
3. 영업경비: 연간관리비의 75\%(감가상각비 미포함) \par
4. 감정평가사甲은 최근 부동산시장의 동향을 분석한 결과 오피스 대비 상가의 수요가 하락추세이고, 배달문화와 비대면문화의 확산으로 오피스보다 상가의 대출금리가 높게 책정된 상황임. 이러한 상황을 고려하여 시장추출법으로 종합환원율을 결정함 \par
{\small\centering\renewcommand{\arraystretch}{1.25}
\resizebox{\ifdim\width>\linewidth\linewidth\else\width\fi}{!}{
\begin{tabular}{|c|c|c|c|}
\hline
\makecell{구분} & \makecell{사례1} & \makecell{사례2} & \makecell{사례3} \\
\hline
\makecell{사례 집합건물} & \makecell{X빌딩 101호} & \makecell{Y빌딩 301호} & \makecell{Z빌딩 101호} \\
\hline
\makecell{용도} & \makecell{근린생활시설} & \makecell{업무시설} & \makecell{근린생활시설} \\
\hline
\makecell{계약면적(원/m\textsuperscript{2})} & \makecell{1,500} & \makecell{1,800} & \makecell{1,200} \\
\hline
\makecell{연간 순영업소득(NOI)(원)} & \makecell{750,000,000} & \makecell{600,000,000} & \makecell{660,000,000} \\
\hline
\makecell{거래가격(원)} & \makecell{15,000,000,000} & \makecell{12,000,000,000} & \makecell{11,000,000,000} \\
\hline
\makecell{거래시점} & \makecell{2020.03.15} & \makecell{2020.03.25} & \makecell{2020.03.30} \\
\hline
\end{tabular}}\par}

※ 사례는 대상 부동산의 인근지역 내 소재함 \par
5. 보증금운용이율: 연 2.0\% \par
< 자료 6 > 영업권평가 참고자료 \par
1. 재무상태표 \par
{\small\centering\renewcommand{\arraystretch}{1.25}
\resizebox{\ifdim\width>\linewidth\linewidth\else\width\fi}{!}{
\begin{tabular}{|c|c|c|}
\hline
\multirow{2}{*}{\makecell{과목}} & \makecell{2019.12.31 현재} & \makecell{2020.09.19 현재} \\
\cline{2-3}
 & \makecell{금액(원)} & \makecell{금액(원)} \\
\hline
\makecell{자산} &  &  \\
\hline
\makecell{Ⅰ.유동자산} & \makecell{40,000,000,000} & \makecell{35,000,000,000} \\
\hline
\makecell{Ⅱ.비유동자산} &  &  \\
\hline
\makecell{1.투자자산} & \makecell{22,000,000,000} & \makecell{21,500,000,000} \\
\hline
\makecell{2.유형자산} & \makecell{20,000,000,000} & \makecell{19,500,000,000} \\
\hline
\makecell{자산총계} & \makecell{82,000,000,000} & \makecell{76,000,000,000} \\
\hline
\makecell{부채} &  &  \\
\hline
\makecell{Ⅰ.유동부채} &  &  \\
\hline
\makecell{1.외상매입금} & \makecell{20,000,000,000} & \makecell{20,000,000,000} \\
\hline
\makecell{2.단기차입금} & \makecell{5,000,000,000} & \makecell{0} \\
\hline
\makecell{Ⅱ.비유동부채} &  &  \\
\hline
\makecell{1. 장기차입금} & \makecell{25,000,000,000} & \makecell{25,000,000,000} \\
\hline
\makecell{부채총계} & \makecell{50,000,000,000} & \makecell{45,000,000,000} \\
\hline
\makecell{자본총계} & \makecell{32,000,000,000} & \makecell{31,000,000,000} \\
\hline
\makecell{부채 및 자본총계} & \makecell{82,000,000,000} & \makecell{76,000,000,000} \\
\hline
\end{tabular}}\par}

※ 2019년 12월 31일 기준 재무상태표상 자산, 부채, 자본의 규모는 과거 5개년 평균과 유사한 수준임 \par
※ 2020년 9월 19일 기준 재무상태표는 비유동자산을 제외하고 공정가치로 조정되었음 \par
※ 2020년 9월 15일 단기차입금(5,000,000,000원)을 현금으로 상환하였으나, 회사의 통상적인 영업경비 충당을 위해 2020년 9월 30일에 재차입예정임 \par
2. 비영업용자산이 포함된 주식회사B의 기업가치는 70,000,000,000원으로 평가하였음 \par
3. 주식회사B 소유 부동산 중 자가사용 외의 집합건물은 임대수익을 얻기 위해 보유하고 있음 \par
4. 비유동자산은 주식회사B 소유 부동산 외에 다른 자산은 없다고 가정함 \par
5. 대상 기업의 무형자산은 영업권만 존재한다고 가정함 \par
\end{longquestion}
